\chapter{Geltungsbereich und Grundbegriffe}

Dieser Leitfaden ergänzt Band~11 um eine ausgewählte Übersicht
zentraler Definitionen und Sätze über \emph{Äquivalenzrelationen}.
Er enthält keine Beweise oder Beweisskizzen. Grundlage ist die
Mengenlehre in ZF ohne zusätzliches Auswahlaxiom; die untersuchten
Relationen tragen ausschließlich die drei Äquivalenzaxiome.
Die Aussagen gelten auf beliebigen Mengen, einschließlich der
leeren Menge, soweit keine Einschränkung genannt ist.

\section{Gegenstand und Abgrenzung}

Der Leitfaden behandelt die Charakterisierung von Äquivalenzrelationen,
ihre Klassen und Partitionen, gesättigte Teilmengen, Konstruktionen,
erzeugte Äquivalenzen und wiederholte Quotientenbildung.
Mengenoperationen dienen als vorausgesetzte Werkzeuge.
Einzelne Sätze nennen zusätzliche Bedingungen, etwa die Inklusion
zweier Äquivalenzrelationen; daraus werden keine weiteren
axiomatischen Strukturen entwickelt.

Allgemeine Funktionentheorie sowie Injektivität, Surjektivität
und Bijektivität gehören zu den Bänden~5--8. Daher entfallen hier
Quotientenprojektionen, Funktionsabstieg, universelle
Abbildungseigenschaften und Gleichmächtigkeit.
Ordnungsrelationen und Verbände, algebraische Kongruenzen,
Topologie, Kardinalitätsformeln und die Auswahl von
Repräsentantensystemen bleiben ebenfalls außerhalb des Leitfadens.
Die Quotientenmenge selbst gehört dagegen zum Kern der
Äquivalenztheorie.

Die Zusammenstellung ist eine Referenzliste wichtiger bekannter
Resultate, keine vollständige Aufzählung aller bekannten oder
ableitbaren Aussagen. Die Nummern mit Präfix $11H$ gehören
ausschließlich zu diesem Leitfaden.

\section{Äquivalenzrelation, Klassen und Partitionen}

\begin{definition}[Äquivalenzrelation]\label{b11guide:def:equivalence}
Sei $A$ eine Menge. Eine Teilmenge $E\subseteq A\times A$
heißt \emph{Äquivalenzrelation auf $A$}, wenn für alle $x,y,z\in A$
gilt:
\[
\begin{aligned}
 &x\sim_E x &&\text{(Reflexivität)},\\
 &x\sim_E y\ \Rightarrow\ y\sim_E x &&\text{(Symmetrie)},\\
 &(x\sim_E y\ \land\ y\sim_E z)\ \Rightarrow\ x\sim_E z
 &&\text{(Transitivität)}.
\end{aligned}
\]
Dabei bedeutet $x\sim_E y$ stets $(x,y)\in E$.
Der Buchstabe $E$ bezeichnet die Paarmenge der Relation,
$\sim_E$ ihren Relationsoperator in Infixschreibweise.
Die Menge aller Äquivalenzrelationen auf $A$ sei
\[
 \EqSet(A)=\{E\in\mathcal P(A\times A)\mid
                       E\text{ ist eine Äquivalenzrelation auf }A\}.
\]
\end{definition}

\begin{definition}[Äquivalenzklasse und Quotientenmenge]
\label{b11guide:def:class}
Für $E\in\EqSet(A)$ und $x\in A$ heißen
\[
 [x]_E=\{y\in A\mid x\sim_E y\},
 \qquad A/E=\{[x]_E\mid x\in A\}
\]
die \emph{Äquivalenzklasse von $x$} und die \emph{Quotientenmenge}.
Ein $x$ mit $C=[x]_E$ heißt \emph{Repräsentant} der Klasse $C$.
Diese Bezeichnung trifft keine gleichzeitige Auswahl für alle Klassen.
\end{definition}

\begin{definition}[Partition]\label{b11guide:def:partition}
Eine \emph{Partition} oder \emph{Zerlegung} von $A$ ist eine
Mengenfamilie $\mathcal C\subseteq\mathcal P(A)$ mit
\[
 \varnothing\notin\mathcal C,\qquad
 \bigcup\mathcal C=A,\qquad
 C,D\in\mathcal C,\ C\ne D\ \Rightarrow\ C\cap D=\varnothing.
\]
Ihre Elemente heißen \emph{Blöcke}. Die leere Familie ist eine
Partition von $\varnothing$; $\{\varnothing\}$ ist keine Partition.
\end{definition}

\section{Wegweiser}

Kapitel~\ref{b11guide:ch:classes} enthält Charakterisierungen und
den Zusammenhang zwischen Klassen und Partitionen.
Kapitel~\ref{b11guide:ch:saturation} sammelt die Sätze über
gesättigte Teilmengen.
Kapitel~\ref{b11guide:ch:constructions} behandelt Einschränkungen,
Produkte, Durchschnitte und Ununterscheidbarkeit.
Kapitel~\ref{b11guide:ch:generated} beschreibt erzeugte
Äquivalenzrelationen, Vereinigungen und Verkettungen.
Kapitel~\ref{b11guide:ch:quotients} enthält den Vergleich von
Klassen und die zweimalige Quotientenbildung.

